% !TEX root = ../ApproxDA-TransUNet.tex
\section{Conclusion}
\label{sec:conclusion}
This paper presented ApproxDA-TransUNet, a dual-attention approximation framework for medical image segmentation featuring independently controllable window size, projection rank, and channel grouping. Across three representative clinical benchmarks, ApproxDA-TransUNet achieves consistent Dice gains over DA-TransUNet with appropriate window configurations, demonstrating that attention approximation directly modulates representational capacity and accuracy rather than acting solely as an efficiency mechanism.

Our systematic ablations further reveal that context sensitivity diverges substantially across clinical targets, a property we formalize via empirical Global Context Sensitivity (GCS): multi-organ CT segmentation requires precise scale tuning, whereas polyp and skin lesion targets exhibit broad robustness across receptive window scales. In addition, our empirical findings highlight that learnable dual-branch routing tends toward static uniform blending on multi-organ CT without surpassing single-branch PAM routing.

\textbf{Limitations and Future Work.} In accordance with benchmark literature conventions, reported metrics reflect single-run performance on standard splits, and certain baseline figures are cited from published studies under varying training protocols. Furthermore, while GCS effectively summarizes empirical context sensitivity across the evaluated benchmarks, developing predictive context heuristics prior to exhaustive sweeping and validating them across broader anatomical modalities and architectures remain valuable directions for future study. Finally, translating theoretical attention-matrix reductions into end-to-end inference acceleration via specialized kernel implementations and lightweight backbones represents a natural next step. We hope this work encourages principled, task-aware exploration of attention approximation across medical imaging domains.
